\section{Introduction}

\subsection{Seventy Years of Convention}

Electroencephalography has recognized discrete frequency bands since Berger described the alpha rhythm in 1929 and subsequent workers named theta, delta, beta, and gamma oscillations across the following decades. The canonical boundaries---delta below $\sim$4 Hz, theta 4--8 Hz, alpha 8--13 Hz, beta 13--30 Hz, gamma above 30 Hz---were established through clinical observation and informal consensus, not derivation from first principles \parencite{buzsaki2006rhythms}. The Greek-letter nomenclature reflects the historical order of discovery, not the frequency order: alpha was observed first, then delta (the slowest), then theta and beta. The resulting naming sequence---delta, theta, alpha, beta, gamma---is neither alphabetical nor frequency-ordered, a small absurdity that underscores the atheoretical character of the enterprise.

Boundary values vary across laboratories, textbooks, and clinical guidelines by several hertz in either direction. The upper limit of theta is variously placed at 7, 8, or 9 Hz; the lower limit of gamma ranges from 25 to 40 Hz. These are not minor calibration disagreements---a 7 Hz versus 9 Hz theta boundary changes which oscillatory peaks are classified as theta versus alpha for a substantial fraction of the population, given that individual alpha frequency (IAF) varies from roughly 8 to 12 Hz across healthy adults \parencite{klimesch1999eeg,corcoran2018iaf,grandy2013iaf}. Data-driven methods confirm that population-optimal boundaries differ from convention and vary across individuals and clinical groups \parencite{cohen2021gedbounds,talebi2022datadriven}. The boundaries are approximate, yet they structure nearly every quantitative EEG analysis performed in the last half century.

The reason the field has tolerated this approximation is that no principled alternative has been established. Several theoretical frameworks have been proposed---examined in detail below---but none has been subjected to formal statistical comparison against competing models on large-scale empirical data. The present study fills this gap.

\subsection{The Log-Frequency Hypothesis}

The most influential theoretical claim about EEG band organization is that band centers are logarithmically spaced---specifically, that the ratio between neighboring band center frequencies approximates a fixed mathematical constant. Three constants have been proposed, and the debate among them has proceeded for over two decades without resolution.

\textcite{penttonen2003natural} compiled known oscillation frequencies from mammalian electrophysiology and observed that they form an arithmetic progression on a natural logarithmic scale, with neighboring bands separated by a factor of approximately Euler's number $e \approx 2.72$. This ``Natural Logarithm Linear Law'' (N3L) was the first formal proposal that inter-band spacing is governed by a specific irrational constant. The theoretical motivation was biophysical: each band is generated by a distinct neural circuit mechanism, and the irrational spacing ratio ensures minimal cross-frequency interference by preventing phase locking between neighboring bands. \textcite{buzsaki2004neuronal,buzsaki2006rhythms,buzsaki2013scaling} elaborated this into a hierarchical oscillatory framework in which the phase of slower oscillations modulates the amplitude of faster ones, creating a ``neural syntax'' across timescales. The logarithmic spacing is, in this view, a design principle of the nervous system preserved across mammalian evolution despite thousand-fold variation in brain size.

An alternative constant---the golden ratio $\phisym = (1 + \sqrt{5})/2 \approx 1.618$---was proposed by \textcite{pletzer2010golden}, who showed mathematically that $\phisym$-related oscillators achieve maximal desynchronization: because $\phisym$ is the ``most irrational'' number in a precise number-theoretic sense, oscillators at $\phisym$-related frequencies can never synchronize their excitatory phases. Analyzing resting EEG, they observed that spectral peaks in standard bands approximate a geometric $\phisym$-series with centers near 3, 5, 8, 13, 21, and 34 Hz---resembling the Fibonacci sequence. A biophysical mechanism was provided by \textcite{roopun2008period}, who demonstrated in rat neocortical slices that gamma ($\sim$40 Hz) and beta2 ($\sim$25 Hz) oscillations have a frequency ratio approximating $\phisym$, and that their period concatenation (adding successive cycle lengths) generates beta1 ($\sim$15 Hz), producing the Fibonacci recurrence $f_{k+2} = f_{k+1} + f_k$ from concrete circuit dynamics. \textcite{kramer2023golden} formalized this into a ``golden rhythms'' framework in which the complete sequence $f_k = \fzero \times \phisym^k$ uniquely satisfies both the segregation requirement (maximal desynchronization) and an integration requirement (cross-frequency coupling when needed), and provided computational support from coupled oscillator simulations.

\textcite{klimesch2018frequency} proposed a hybrid coordinate system anchoring band \emph{centers} at factor-2 intervals (2.5, 5, 10, 20, 40 Hz) while defining band \emph{boundaries} using $\phisym$, yielding non-overlapping bands with lower limits at $f_{d(i-1)} \times \phisym$ and upper limits at $f_{d(i+1)} / \phisym$. This dual-constant system is motivated by the distinct functional roles of the two constants: factor-2 (octave) spacing supports stable phase-phase coupling between hierarchically related oscillations, while $\phisym$-defined boundaries minimize cross-talk. Klimesch's framework is notable for extending below 1 Hz to encompass heart rate, respiration, and blood pressure oscillations, treating brain and body as a unified frequency architecture.

Most recently, \textcite{ursachi2026phi} validated $\phisym$-organization as an individual differences variable in 320 subjects across two independent datasets, introducing the Phi Coupling Index (PCI) to quantify the proximity of each individual's theta--alpha centroid ratio to $\phisym$. In a companion preprint, \textcite{ursachi2026boundaries} reported that three independent data-driven boundary-detection methods converge on boundaries whose successive ratios approximate $e - 1 \approx 1.718$ rather than $\phisym$, and demonstrated through power analysis that studies with $N < 30$ subjects systematically favor $\phisym$ over $e-1$ due to insufficient statistical power.

The three candidate constants---$e \approx 2.718$, $\phisym \approx 1.618$, and $e - 1 \approx 1.718$---make different empirical predictions. They are not mutually exclusive in principle (centers and boundaries could obey different constants, as in Klimesch's framework), but they have never been compared against each other, or against non-log alternatives, on the same dataset with formal model selection. The Penttonen \& Buzsaki observation that launched this field was a descriptive compilation of known band centers; it included no null hypothesis test, no confidence intervals on the ratio estimate, and no comparison to $\phisym$, $e-1$, or linear spacing. Twenty-three years of theoretical elaboration rest on this foundation.

\subsection{Within-Band Organization as an Individual Difference}

A separate problem concerns what happens \emph{inside} each band. Most theoretical frameworks address inter-band spacing---where the boundaries fall---but say little about how oscillatory peaks distribute within a band, or whether that distribution varies meaningfully across individuals.

The dominant EEG individual-differences measures each capture one aspect of spectral organization. Band power indexes total oscillatory energy but conflates periodic and aperiodic components \parencite{donoghue2020fooof}. IAF identifies the alpha peak location and predicts cognitive performance, age, and clinical status, but describes only a single band and a single feature of its spectrum \parencite{klimesch1999eeg,corcoran2018iaf,grandy2013iaf}. The aperiodic slope of the $1/f$ background reflects excitation--inhibition balance and changes with development, aging, and psychopathology \parencite{voytek2015aperiodic}, but by definition captures activity distributed uniformly across all frequencies, not band-specific structure. Connectivity and cross-frequency coupling measures \parencite{keitel2016spectral,mahjoory2020frequency} characterize relationships between regions and between frequency scales, but not the within-band frequency structure at a single recording site.

None of these measures captures whether, within a given band, oscillatory peaks cluster at a preferred frequency or spread across the band. A brain whose theta activity concentrates near 6 Hz has a very different theta profile from one whose theta activity spreads evenly from 4 to 8 Hz---yet these two patterns have identical band power, similar aperiodic slopes, and comparable connectivity profiles. The shape of the within-band peak distribution is, to our knowledge, unmeasured.

We term this shape \textbf{spectral differentiation}: the degree to which a brain's oscillatory peaks concentrate at specific within-band positions rather than distributing uniformly. A highly differentiated within-band spectrum has steep enrichment gradients---peaks cluster at preferred positions, with voids elsewhere. A flat, undifferentiated spectrum has peaks distributed without preference across the band. This construct is operationally distinct from power (it is normalized), from IAF (it captures band shape, not peak location), and from connectivity (it does not require multi-electrode data). Whether it captures biologically meaningful variance is an empirical question that the present study addresses.

\subsection{The Present Study}

We report three results. First, we provide the first formal statistical demonstration that EEG oscillatory peaks are logarithmically organized, using model comparison across seven candidate scaling ratios plus linear spacing on 4,572,489 FOOOF-detected peaks from 9 independent datasets spanning 2,097 subjects aged 5--70. We show that log-frequency description is decisively superior to linear ($\Delta$BIC $> 50$), that band boundaries are genuine spectral features rather than $1/f$ artifacts (aperiodic null: $p < 0.0001$ at all boundaries), and that among fixed-ratio models, $\phisym$ best fits empirical boundary positions---while rejecting $e$ as the scaling constant proposed by \textcite{penttonen2003natural}.

Second, we optimize the EEG frequency coordinate system by sweeping a two-dimensional parameter space of base frequency \fzero\ and scaling ratio, identifying which combination produces the simplest, most consistent within-band enrichment profiles. The $\phisym$-lattice---bands anchored to \fzero\ and separated at \fzero\ $\times \phisym^n$---achieves profile simplicity $R^2 = 0.968$ and places all four inter-band boundaries within 2\% of their empirically optimal positions. Within bands, however, no fine-structure exists at $\phisym$-specific positions: enrichment profiles are smooth ramps or mountains with no preferred landmarks, and noble-number positions carry no more signal than simple rationals ($p = 0.67$--$0.98$). The value of the $\phisym$-lattice is therefore its boundary placement, not any hypothesized computation with irrational numbers.

Third, we establish spectral differentiation as a novel EEG biomarker. Using the $\phisym$-lattice coordinate system to define within-band positions, we compute per-subject enrichment profiles and derive interpretable scalar metrics (band mountain, U-shape, ramp depth, center depletion). These metrics predict cognitive performance in the LEMON dataset (31 FDR-significant correlations across four frequency bands and four cognitive tests), track an inverted-U developmental trajectory from childhood through young adulthood in HBN ($N = 927$, 60 FDR-significant age correlations) followed by decline in aging in the Dortmund Vital Study ($N = 608$, 41 FDR-significant correlations), dissociate externalizing from internalizing psychopathology in HBN (18 and 7 FDR survivors respectively), and demonstrate 5-year test--retest reliability (ICC $= 0.75$ for the most reliable metric). Spectral differentiation is uninformative about personality traits (0 of 11,970 tests survive FDR), medical and metabolic variables (0 FDR), and handedness (0 FDR), suggesting domain specificity rather than a generic signal artifact.

This study also corrects three systematic biases in a prior report \parencite{lacy2026frontiers} that introduced the $\phisym$-lattice framework: a mismatch between the extraction base frequency ($\fzero = 7.83$ Hz) and the lattice base frequency ($\fzero = 7.60$ Hz), a linear-Hz weighting of the null hypothesis that inflated enrichment estimates at high frequencies, and separate FOOOF fitting of theta and alpha bands that created an artificial boundary artifact. Correcting these biases substantially revises the enrichment landscape, reverses the sign of several previously reported effects, and provides a more reliable foundation for the biomarker analyses.

Not all of these findings carry equal evidential weight. The frequency-scaling and boundary results are robust across all datasets, methods, and analytical choices. The biomarker results---developmental trajectory, test--retest reliability, and the strongest cognitive and psychopathology associations---replicate across parametric (FOOOF) and non-parametric (IRASA) extraction methods, but the breadth of the cognitive map is FOOOF-specific and relies on a single adult sample. A mechanistic interpretation linking spectral differentiation to GABAergic inhibitory circuits is offered as a hypothesis-generating framework, not a demonstrated result. The Discussion distinguishes these tiers explicitly.
